\documentclass[12pt,a4paper]{article}
\usepackage[utf8]{inputenc}
\usepackage[vietnamese]{babel}
\usepackage{graphicx}
\usepackage{geometry}
\usepackage{hyperref}
\geometry{a4paper, margin=1in}

\title{\textbf{BÁO CÁO DỰ ÁN CƠ SỞ DỮ LIỆU (DBI202)}\\
\large Hệ thống Quản lý Ca làm việc và Chấm công (Staff Shift Management)}
\author{Nhóm Sinh Viên}
\date{\today}

\begin{document}

\maketitle

\begin{abstract}
Báo cáo này trình bày thiết kế và triển khai một hệ thống cơ sở dữ liệu để quản lý ca làm việc và chấm công của nhân viên. Dự án bao gồm việc phân tích yêu cầu, xây dựng Mô hình Thực thể - Liên kết (ERD), thiết kế cơ sở dữ liệu quan hệ trên SQL Server và sử dụng ngôn ngữ JavaScript (Node.js) để thu thập, xử lý dữ liệu.
\end{abstract}

\tableofcontents
\newpage

\section{Giới thiệu (Introduction)}
Trong môi trường doanh nghiệp hiện đại, việc quản lý thời gian làm việc của nhân viên một cách chính xác là cực kỳ quan trọng. Dự án \textbf{Staff Shift Management} được xây dựng nhằm mục đích tin học hóa quy trình phân ca và chấm công. Hệ thống cho phép:
\begin{itemize}
    \item Lưu trữ thông tin nhân sự và các ca làm việc.
    \item Lên lịch phân ca cho từng nhân viên (quan hệ nhiều - nhiều).
    \item Ghi nhận thẻ chấm công (Check-in, Check-out) hàng ngày.
    \item Xử lý dữ liệu thời gian thực để tính toán giờ làm và đánh giá trạng thái đi làm (đúng giờ, đi trễ).
\end{itemize}

\section{Mô hình Dữ liệu (Data Model)}

\subsection{Mô hình Thực thể - Liên kết (ERD)}
Dựa trên phân tích yêu cầu, hệ thống bao gồm 3 thực thể chính:
\begin{enumerate}
    \item \textbf{STAFF (Nhân viên):} Quản lý thông tin cá nhân và mức lương theo giờ.
    \item \textbf{SHIFT (Ca làm việc):} Quản lý khung giờ làm việc.
    \item \textbf{TIMECARD (Thẻ chấm công):} Lưu trữ giờ vào, giờ ra của nhân viên theo ngày.
\end{enumerate}

\noindent Các mối quan hệ trong mô hình:
\begin{itemize}
    \item \textbf{SCHEDULES:} Mối quan hệ nhiều-nhiều (M-N) giữa STAFF và SHIFT.
    \item \textbf{LOGS \& SPECIFIES:} Mối quan hệ một-nhiều (1-N) nối từ STAFF và SHIFT đến TIMECARD.
\end{itemize}
\textit{(Ghi chú: Tại Overleaf, nhóm có thể tải file ảnh ERD.png lên và chèn vào đây bằng lệnh includegraphics)}

\subsection{Thiết kế Cơ sở dữ liệu (Database Schema)}
Hệ quản trị CSDL được lựa chọn là Microsoft SQL Server. Cấu trúc các bảng như sau:
\begin{verbatim}
-- Bảng STAFF
CREATE TABLE STAFF (
    UserID INT PRIMARY KEY, FullName NVARCHAR(100), 
    Role NVARCHAR(50), HourlyRate DECIMAL(10, 2)
);

-- Bảng SHIFT
CREATE TABLE SHIFT (
    ShiftID INT PRIMARY KEY, ShiftName NVARCHAR(50), Time_Range NVARCHAR(50)
);

-- Bảng TIMECARD
CREATE TABLE TIMECARD (
    TimecardID INT PRIMARY KEY, WorkDate DATE, 
    CheckIn TIME, CheckOut TIME,
    UserID INT FOREIGN KEY REFERENCES STAFF(UserID),
    ShiftID INT FOREIGN KEY REFERENCES SHIFT(ShiftID)
);
\end{verbatim}

\section{Quy trình Thu thập và Xử lý Dữ liệu (Data Workflow)}
Tuân thủ yêu cầu học phần tuần 6, nhóm thiết lập một luồng xử lý dữ liệu như sau:
\subsection{Thu thập (Data Collection)}
Dữ liệu thô được thu thập từ hai nguồn chính:
\begin{itemize}
    \item \textbf{Dữ liệu tĩnh:} Thông tin Nhân viên và Ca làm việc do bộ phận Quản lý (HR) nhập.
    \item \textbf{Dữ liệu động:} Thẻ chấm công (CheckIn/CheckOut) được lấy từ thiết bị phần cứng.
\end{itemize}

\subsection{Xử lý dữ liệu với Node.js (Data Processing)}
Nhóm sử dụng \texttt{JavaScript (Node.js)} để thực hiện:
\begin{itemize}
    \item \textbf{Tính toán giờ làm (TotalHours):} Xử lý các phép toán thời gian cơ bản, xử lý các trường hợp ca làm qua đêm (âm giờ).
    \item \textbf{Phân tích trạng thái (Status):} Dùng logic điều kiện (If-Else) để dán nhãn "Đi trễ" hay "Đúng giờ" dựa trên mốc thời gian quy định của từng ca làm.
\end{itemize}

\section{Kết luận (Conclusion)}
Hệ thống \textbf{Staff Shift Management} đã đáp ứng tốt các yêu cầu cơ bản về quản trị cơ sở dữ liệu. Thông qua dự án này, nhóm đã nắm vững kiến thức về thiết kế ERD, triển khai mô hình quan hệ trên SQL Server và đặc biệt là tích hợp Node.js để xử lý, làm sạch dữ liệu thực tế.

\end{document}
